\begin{figure}[htbp]
    \centering
    \resizebox{0.92\linewidth}{!}{%
    \begin{tikzpicture}[
        font=\sffamily,
        corpusBox/.style={
            draw, thick, align=center,
            inner xsep=6pt, inner ysep=6pt,
            text width=2.3cm, minimum height=1.3cm
        },
        unifiedBox/.style={
            draw, thick, align=center,
            inner xsep=10pt, inner ysep=8pt,
            text width=5.0cm, minimum height=1.5cm,
            fill=blue!5
        },
        modelBox/.style={
            draw, thick, align=center,
            inner xsep=8pt, inner ysep=8pt,
            text width=3.9cm, minimum height=1.6cm
        },
        finalBox/.style={
            draw, thick, align=center,
            inner xsep=8pt, inner ysep=8pt,
            text width=4.2cm, minimum height=1.8cm,
            fill=blue!10
        },
        arrow/.style={-{Stealth[scale=1]}, thick},
        lbl/.style={align=center, font=\bfseries}
    ]

    % ---------- Fila 1: corpus individuales ----------
    \node[corpusBox] (pharmaconer) at (0,0)
        {\footnotesize\textbf{PharmaCoNER}\\\scriptsize Fármacos y proteínas\\\scriptsize(\textit{SPACCC})};
    \node[corpusBox, right=0.45cm of pharmaconer] (symptemist)
        {\footnotesize\textbf{SympTEMIST}\\\scriptsize Síntomas};
    \node[corpusBox, right=0.45cm of symptemist] (medprocner)
        {\footnotesize\textbf{MedProcNER}\\\scriptsize Procedimientos};
    \node[corpusBox, right=0.45cm of medprocner] (distemist)
        {\footnotesize\textbf{DisTEMIST}\\\scriptsize Enfermedades};

    % Punto medio de toda la fila 1, para centrar el título
    \coordinate (rowMid) at ($(pharmaconer.north)!0.5!(distemist.north)$);
    \node[lbl, above=0.3cm of rowMid] (corpusTitle) {Campañas de evaluación del BSC};

    % Punto medio entre SympTEMIST y MedProcNER, para centrar el corpus unificado
    \coordinate (midTop) at ($(symptemist.east)!0.5!(medprocner.west)$);

    % ---------- Fila 2: corpus unificado ----------
    \node[unifiedBox, below=0.9cm of midTop, anchor=north] (unified) {
        \textbf{Corpus Multiclase Unificado}\\[3pt]
        \footnotesize Reconocimiento y vinculación (\textit{entity linking}) de entidades biomédicas en español, mapeadas a \textit{SNOMED CT}
    };

    % ---------- Fila 3: modelos derivados (simétricos bajo unified) ----------
    \node[modelBox, below=0.9cm of unified, anchor=north, xshift=-2.4cm] (robertaModel) {
        \footnotesize\textbf{RoBERTa-base-biomedical-clinical-es}\\[3pt]
        \scriptsize Preentrenado con literatura médica y registros de salud electrónicos
    };

    \node[finalBox, below=0.9cm of unified, anchor=north, xshift=2.5cm] (systemModels) {
        \footnotesize\textbf{Modelos empleados en el sistema}\\[3pt]
        \footnotesize\textbf{spanish-medical-ner}\\
        \scriptsize Síntomas, procedimientos y enfermedades\\[3pt]
        \footnotesize\textbf{bsc-bio-ehr-es-pharmaconer}\\
        \scriptsize Fármacos y proteínas
    };

    % ---------- Conexiones: corpus individuales -> unificado ----------
    \draw[arrow] (pharmaconer.south) |- (unified.north -| pharmaconer);
    \draw[arrow] (symptemist.south) -- (unified.north -| symptemist);
    \draw[arrow] (medprocner.south) -- (unified.north -| medprocner);
    \draw[arrow] (distemist.south) |- (unified.north -| distemist);

    % ---------- Conexiones: unificado -> modelos derivados ----------
    \draw[arrow] (unified.south) -- (robertaModel.north)
        node[font=\scriptsize, pos=0.45, left=2pt] {entrena};
    \draw[arrow] (unified.south) -- (systemModels.north)
        node[font=\scriptsize, align=center, pos=0.45, right=2pt] {hereda criterios\\de anotación};

    % ---------- Conexión: RoBERTa -> modelos del sistema ----------
    \draw[arrow] (robertaModel.east) -- (systemModels.west)
        node[font=\scriptsize, midway, above] {linaje};

    \end{tikzpicture}%
    }
    \caption{Ecosistema de corpus y modelos biomédicos en español desarrollados por el Barcelona Supercomputing Center (BSC), y su relación con los modelos empleados en el sistema propuesto.}
    \label{fig:diagrama_ecosistema_bsc}
\end{figure}